% Created 2018-10-22 Mon 15:35
% Intended LaTeX compiler: pdflatex
\documentclass[11pt]{article}
\usepackage[utf8]{inputenc}
\usepackage[T1]{fontenc}
\usepackage{graphicx}
\usepackage{grffile}
\usepackage{longtable}
\usepackage{wrapfig}
\usepackage{rotating}
\usepackage[normalem]{ulem}
\usepackage{amsmath}
\usepackage{textcomp}
\usepackage{amssymb}
\usepackage{capt-of}
\usepackage{hyperref}
\pagenumbering{roman}
\usepackage[margin=0.80in]{geometry}
\usepackage[backend=biber, style=ieee, url=false]{biblatex}
\usepackage{float}
\usepackage[super,negative]{nth}
\usepackage[capitalise]{cleveref}
\usepackage{pst-node,transparent,ragged2e}
\addbibresource{/home/nlisa/.spacemacs.d/org-files/bibliography.bib}
\DeclareFieldFormat[inproceedings]{citetitle}{\textit{#1}}
\DeclareFieldFormat[inproceedings]{title}{\textit{#1}}
\DeclareFieldFormat[misc]{citetitle}{#1}
\DeclareFieldFormat[misc]{title}{#1}
\renewcommand*{\bibpagespunct}{%
\ifentrytype{inproceedings}
{\addspace}
{\addcomma\space}}
\AtEveryCitekey{\ifuseauthor{}{\clearname{author}}}
\AtEveryBibitem{\ifuseauthor{}{\clearname{author}}}
\SpecialCoor
\author{Nyameko Lisa}
\date{19 October 2018}
\title{IPSP089 - Legal Aspects of Traditional Knowledge and Biodiversity\\\medskip
\large Portfolio Exam: 743336}
\hypersetup{
 pdfauthor={Nyameko Lisa},
 pdftitle={IPSP089 - Legal Aspects of Traditional Knowledge and Biodiversity},
 pdfkeywords={},
 pdfsubject={},
 pdfcreator={Emacs 26.1 (Org mode 9.1.14)},
 pdflang={English}}
\begin{document}

\maketitle
\addvspace{110pt}
\centering{
\pnode(0.5\textwidth,-0.5\textheight){thisCenter}
\rput(thisCenter){%\transparent{0.25}
\includegraphics[width=2.7in]{/home/nlisa/course/llb/wipo-unisa/UNISACoatofArms.eps}}}

\justifying
\addvspace{110pt}
\section*{Declaration}
\label{sec:org2a03ac6}
I know that plagiarism is to use someone else’s work and pass it off as my own.\\
I know that plagiarism is wrong.\\
I confirm that this assignment is my own work.\\
I have acknowledged in the bibliography accompanying the assignment all the sources that I have used.\\
I have not directly copied without acknowledgement anything from the Internet or from any other source.\\
I have indicated every quotation and citation in a footnote or bracket linked to that quotation.\\
I have not allowed anyone else to copy my work and to pass it off as their own work.\\
I understand that if any unacknowledged copying whatsoever appears in my assignment I will receive zero per cent for the assignment.\\
I am aware of the UNISA policy on plagiarism and understand that disciplinary proceedings can be instituted against me by UNISA if I contravene this policy.\\
I indicate my understanding and acceptance of this declaration by
entering my name here under:
\begin{itemize}
\item Name: \textbf{Nyameko Lisa} (Student Number: \textbf{7874-909-3})
\end{itemize}

\subsection*{NOTE}
\label{sec:org9d1fc7d}
Please note that footnotes will be denoted as \footnote{This is a footnote.} and will
appear at the bottom of the page.\\
References will be denoted by \cite{geyer10_towards_clearer_defin_underst_indi} and will appear at the end of the document.
\newpage
\pagenumbering{arabic}

\section{Critically analyse this statement:}
\label{sec:org24862d9}

\textbf{The concept `traditional knowledge' has many facets and encompasses several components. What characterizes traditional knowledge is the fact that, generally, it is not produced systematically, but in accordance with the individual creator's, or the collective creators' responses to their cultural environment and their interaction with it. [25]}\\

Traditional knowledge systems, creations, innovations, or cultural expressions
refer to \emph{literary, artistic or scientific works, performances, scientific
inventions or discoveries, designs, marks or symbols}, that have evolved in
response to a changing environment and are inherited from one generation to
another within a particular group of people or their territory,
\cite{geyer10_towards_clearer_defin_underst_indi}.\\

Informal regimes relied on mechanisms based on non-systematic and undocumented
rules, obligations and rights, that are enforced either by elders or religious
leaders. These regimes are `self-legitimized'
\cite{amechi15_leverag_trad_knowl_medicinal}, within and across the community
though the accrual of consensus to abide by and adhere to these laws and
rules.\\

Generally speaking customary law is not recognised by legal institutions outside
of the relevant communities where they were originally established. The culture
and history of a people are preserved by the community elders and leaders, who
are also responsible for determining eligibility and allocation of various arts
to be practiced amongst the community's artisans, as well as dispute
resolution, \cite{geyer10_towards_clearer_defin_underst_indi}.\\

Many existing legal frameworks are inadequate at resolving or determining a
common denominator between the economic interests of commercialisation, and the
traditional interests of indigenous peoples. Indigenous knowledge is both
vulnerable to misappropriation under existing industrial property systems and
subsequent financial exploitation by global drug industries,
\cite{masango10_indigenous_trad_knowl_protect_prospects}. All without either
consent from, or acknowledgement of the indigenous peoples, who \emph{arguably} could
be considers the original inventors.\\

Globalisation has given rise to a reluctance of younger generations to learn the
traditional and cultural practices, favouring instead the economic interests of
industrial endeavors. Coupled with the absence of willing heirs and
participants, means that there is significant risk of loss of the traditional
and indigenous knowledge should it be inadequately documented and ill-preserved.
This problem also manifests itself through the forced displacement of indigenous
traditional communities based on Eurocentric conservation and biodiversity
management strategies \cite{simelane09_african_trad_knowl_sys_biodiv_manag}.\\

Modern society generally holds a prejudiced view towards traditional and
indigenous knowledge, denigrating it as `primitive'. This lack of appreciation
for such knowledge, stems from the indigenous people traditionally exploiting
the \emph{`effects'} of their traditional knowledge, rather than describing them in
terms of the scientific aspects of the molecular and chemical constituents of
modern biochemical interactions, \cite{amechi15_leverag_trad_knowl_medicinal}.
For example with respect to patent registration as per the provisions of
[section 25(1)] \cite{rsa78_patents_act}, pharmaceutical companies exploit the
\textbf{`inventive step'} requirement for patents by applying it to their method of
extraction and isolation of the active chemical components, that give rise to
said effects.

\section{Discuss how traditional crafts are protected.}
\label{sec:org527e54b}

\textbf{In your answer, refer to relevant case studies, also the use of alternative
forms of intellectual property in dealing with this issue. [25]}\\

As per the provisions of [section 1]\cite{rsa78_copyrightact} traditional crafts
constitute \textbf{`artistic works'} and as such are eligible to enjoy \textbf{copyright
protection} as per the provisions of [section 2(1)(c)]\cite{rsa78_copyrightact}
provided that they meet the \emph{`inherent'} and \emph{`formal'} or \emph{`statutory'}
requirements for copyright to subsist in that work:
\begin{itemize}
\item Inherent Requirements
\begin{itemize}
\item The requirement of material embodiment is clearly satisfied as the
work exists in material form.
\item As per [article 9(2)]\cite{wto17_trips} and [article
2]\cite{wipo96_copyright_treaty}, generally speaking it can be demonstrated
that traditional crafts are a \emph{`particular form of expression of thought'},
thus satisfying the requirement of originality.
\end{itemize}
\item Formal Requirements
\begin{itemize}
\item The author(s) of the traditional craft must be a of, or domiciled in, or a
permanent resident of a Berne Convention country [section 3, section
37]\cite{wipo86_berne}, meaning that they were a \textbf{qualified person} at the time
of creation of the work.
\item Moreover it is implied that the work was first made in the respective Berne
member state, by the corresponding indigenous peoples, in the South African
context [section 4]\cite{rsa78_copyrightact} applies, hence the indigenous
people may allege that the traditional craft meets the statutory
requirements for the subsistence of copyright in the work.
\end{itemize}
\end{itemize}

In the Australian matter \cite{aus95_milpurrurru_v_indofurn} concerning the
importation by a Perth based company of carpets manufactured in Vietnam, that
reproduced Aboriginal designs. While the court awarded substantial damages for
copyright infringement to the Aboriginal artists whose designs had been
reproduced on the carpets, it was unable to compensate the indigenous community
whose images were used in culturally inappropriate ways, where the ruling noted
that copyright legislation is inadequate in recognizing and compensating for the
communal harm resulting from unauthorized reproduction of indigenous designs:
\begin{quote}
\textit{``the statutory remedies do not recognize the infringement of ownership rights of the kind which reside under Aboriginal law in the traditional owners of the dreaming stories'' (at p. 39081)}
\end{quote}
As per the definitions of [section 2(1)]\cite{rsa93_tm_act} and given that the
Republic is a Paris Convention Member State, [article 15(1)]\cite{wto17_trips},
[article 1(2)]\cite{wipo83_paris_conve_protect_ip}, it follows that a traditional
handicraft of an indigenous peoples may be understood to be a \textbf{mark}, in that it
is either a name or sign capable of being represented graphically. Moreover it
constitutes a \textbf{trade mark}, in relation to distinguishing their trade or use
from other goods or services connected in the course of trade regarding a
specific \emph{`class of goods'} as stipulated in \cite{wipo57_ncl}. And as such as per
the provisions of [section 9(1)]\cite{rsa93_tm_act}, it follows that the
traditional craft may qualify for registration as trade mark, subject to
limitations, in relation to use within those limitations, and that none of the
disqualifying conditions are applicable: [paragraphs (1), (4), (5), (8),
(9), (10) or (11) of section 10]\cite{rsa93_tm_act}.\\

As per [section 2(4)]\cite{rsa93_tm_act}, \emph{`use'} or \emph{`proposed use'} shall be
construed as the use of a visual representation of the mark [section
2(2)(a)]\cite{rsa93_tm_act}, as pertaining to physical form or other relation in
terms of goods[section 2(3)(a)]\cite{rsa93_tm_act}, or as pertaining to the use
thereof in any relation to the performance of any services\footnote{Also referred to
as service marks, [article 1(2)]\cite{wipo83_paris_conve_protect_ip}} [section
2(3)(b)]\cite{rsa93_tm_act}. Provided that no others claim \textbf{priority} to the mark
at the time of application for registration, then as per [section
9(2)]\cite{rsa93_tm_act}, the traditional craft mark in questions shall be
considered to be capable of distinguishing within the meaning described above.
Thus also negating the unregistrability conditions of the provisions of
[paragraphs (2), (3) or (6) of section 10]\cite{rsa93_tm_act}.\\

Traditional handicrafts fall under the ambit of either \textbf{`certification trade
marks'} [section 42(1)]\cite{rsa93_tm_act}, [articles 6 and
6\(^{\text{bis}}\)]\cite{wipo83_paris_conve_protect_ip} and [article 15]\cite{wto17_trips} i.e.
those trade marks for which the goods and services are certified by any person
in respect of kind, quality, quantity, intended purpose, value, geographical
origin of the goods or services or their mode or time of production; or
\textbf{`collective trade marks or geographical indications'} [section
43(1)]\cite{rsa93_tm_act}, [articles 6\(^{\text{quinquies}}\) and
7\(^{\text{bis}}\)]\cite{wipo83_paris_conve_protect_ip} and [articles 22 \&
23]\cite{wto17_trips} i.e. those trade marks for which their members are
distinguished by exclusive association. As per the provisions of [section
43(2)]\cite{rsa93_tm_act}, geographical names or other indications of geographical
origin may be registered as collective trade marks.\\

The law of Republic recognizes and grants a general action in the case of
unlawful competition based on the principles of \emph{lex Aquilia.} As per the
quotation from Judge van Dijkhorst in the matter
\citetitle{dijkhorst81_atlas_v_pikkewyn}, \cite{dijkhorst81_atlas_v_pikkewyn}, the
difficulty in establishing an action in the case of unlawful competition is to
quantify the degree of unfairness and/or dishonesty:
\begin{quote}
\textit{``I am not the first nor will I be the last to lament upon the difficulty of determining the dividing line between lawful and unlawful interference with the trade of another…. What is needed is a legal standard firm enough to afford guidance to the Court, yet flexible enough to permit the influence of an inherent sense of fair play''}
\end{quote}
Moreover, as per Judge van Dijkstra, the ideal legal instrument is one in which
the framework for establishing or trying unlawful competition is well defined
and robust enough to provide adequate direction to the courts, but not so rigid
as to restrict the sentiment of an environment of fair play and open
competition.\\

As per [article 10\(^{\text{bis}}\)]\cite{wipo83_paris_conve_protect_ip}, where it is
stipulated that (1) Union member countries must provide adequate protection for
the nationals against unfair competition; (2) dishonest\footnote{Any acts of
competition contrary to honest practices in industrial or commercial matters.}
acts of competition constitutes acts of unfair competition; and finally (3) a
detailed list and description of the forms and nature of acts that would
constitute unlawful competition.\\

In corroboration to the above [article 1]\cite{wipo96_model_provi_unfair_comp}
emphasis es that in the course of industrial or commercial activities an act or
practice that is contrary to honest practices shall be deemed as an act of
unfair competition. Moreover \cite{wipo96_model_provi_unfair_comp} does not
require the entities to be in direct competition, unlike
\cite{wipo83_paris_conve_protect_ip}. The basis for much of the provisions and
notes of \citetitle{wipo96_model_provi_unfair_comp}, are based on those
described in [article 10\(^{\text{bis}}\)]\cite{wipo83_paris_conve_protect_ip}, albeit in much
more detail and specifics.\\

As per [article 40]\cite{wto17_trips}, members of the World Trade Organization,
members are free to determine appropriate methods in implements it's provisions,
moreover regarding competition \cite{wto17_trips} gives particular attention to
intellectual property rights, their abuse and subsequent adverse effects in
relevant markets, trade, transfer and dissemination of technology. Indigenous
peoples may seek relief from a number of forms of unlawful competition, amongst
other grounds on the basis of:
\begin{itemize}
\item \textbf{Passing off}, in light of any similarities of the traditional handicraft with
that of an infringing mark, as per [article
10\(^{\text{bis}}\)(3)(1)]\cite{wipo83_paris_conve_protect_ip}, [article
2(1)]\cite{wipo96_model_provi_unfair_comp} and [article 16(1)]\cite{wto17_trips}.
\item \emph{Damage another's goodwill or reputation}, specifically the \textbf{Dilution of their
goodwill or reputation}, in regard to the lessening of distinctive character
or advertising value of the traditional handicraft, as well as the appearance
and presentation of the product, as per [article
3(2)]\cite{wipo96_model_provi_unfair_comp}
\item \textbf{Unfair competition in respect of secret information}, where another
enterprise has unlawfully acquired the information in the manufacture of
creation of the traditional handicraft, as per [article 39(2)]\cite{wto17_trips}
and [article 6(1)]\cite{wipo96_model_provi_unfair_comp}.
\end{itemize}

As per the definition specified in [section 2]\cite{rsa78_patents_act} and
[article 27(1)]\cite{wto17_trips}, traditional handicrafts need satisfy the provisions for a \textbf{`patent'} to be granted under [section 25]\cite{rsa78_patents_act}, wherein subsection (1) of that section of the Act stipulates that a patent may be granted for a \textbf{new} invention, involving an \textbf{inventive step} and has application to trade or industry.\\

The content of the concept \textbf{`patentable subject-matter'} is usually established
in the negative sense by statue, i.e. inventions explicitly `excluded' from
qualification for the purposes of patent protection. Traditional handicrafts
shall enjoy patent protection provided they neither fall  into any of the
categories listed under [subsections (2) and (3) of section
25]\cite{rsa78_patents_act}, nor those articulated in [articles 27(2) \&
27(3)]\cite{wto17_trips}, nor those of [articles 4 and
5]\cite{wipo83_paris_conve_protect_ip}.\\

Lastly indigenous peoples should also be advised that the legislation makes
provision for both \textbf{`aesthetic designs'} (Part A) as well as \textbf{`functional
designs'} (Part F) \textbf{industrial design} registrations, as per the provisions of
[section 14(1)(a-b)]\cite{rsa93_designs_act}, for the protection of the physical
form of articles of manufacture intended to be multiplied by industrial
processes, [section 14(4)]\cite{rsa93_designs_act}.

\section{Read the following statement and answer the question that follows:}
\label{sec:org7a5b99c}

\textbf{In developed countries, expressions of folklore are generally considered as belonging to the public domain. This approach explains why, in the main, developed countries generally did not establish a legal protection of the manifold national or other communicate interests related to the utilization of expressions of folklore. [25]}

\subsection{Do you agree with the statement? Indicate and explain your reasons. (15)}
\label{sec:orgf06c1c4}

Yes, I concur with the sentiments of that statement. As per the introductory
observations in \cite{wipo85_model_provi_national_laws_folklore}, folklore is an
important element of a nation's cultural heritage, providing its people with an
avenue for self-expression and it exists as a means of a country's social
identity. Traditional expressions have generally been characterized as a broad
set of oral and or written literary, artistic, religious, scientific,
technological works and other traditions and productions that are transmitted
from one generation to another, and whose creators need not necessarily be
indigenous peoples.\\

With regards to folklore, aspects related to its subject matter, technological
and methodological content need to be taken into consideration. Thus it is
important to consider not only protection from a copyright perspective, but also
the mechanisms of industrial property protection, including patents, trademarks
and industrial designs \cite{saurombe09_protection_indigenous_trad_knowl_through_ip}.\\

The need for protection of expressions of folklore arises where the creators or
rights holders, would need to prevent others from reproducing their creations or
disseminating them to the public through performance or broadcast, \emph{(i.e. law of
copyright)}; and to restrict the unauthorized misappropriation of technological
and methodological ideas, inventions, crafts, designs and carvings, \emph{(i.e. laws
of patents, trademarks and industrial designs)}
\cite{masango10_indigenous_trad_knowl_protect_prospects}. Moreover, informal
traditional customary regimes may not provide adequate protection against
cultural, economic or commercial exploitation.\\

Both globalization and the advent of technological telecommunications advances
in recent history, have lead to the improper exploitation of cultural
expressions of folklore currently experienced today \cite{geyer10_towards_clearer_defin_underst_indi}. The misappropriation of the
intellectual property rights that \emph{`should'} be enjoyed by a people, resulting
from their endeavors and accumulated cultural heritage can take a number of
forms, from commercial exploitation to adaptation and distortion of the original
works.\\

There exists a clear need for the development of both national and international
instruments/frameworks, providing legal protection against the exploitation of
cultural expressions of folklore, without any compensation to their originating
communities \cite{saurombe09_protection_indigenous_trad_knowl_through_ip}. Moreover, litigation processes are also required to reclassify
those expressions of folklore that were previously considered to belong to the
\textbf{public domain}. Not forgetting the subsequent infringement and sanctioning
procedures with regards to infringement of the rights in `works of folklore'.

\subsection{Give an indication of what is the position in developing countries. (10)}
\label{sec:orgd4aea81}

A number of developing nations hold a number of standing technical cooperation
programs with WIPO \cite{wto17_trips} and are treaty signatories to the Paris
Convention \cite{wipo83_paris_conve_protect_ip}. As per the provisions of [article
15(4)(a-b)]\cite{wipo86_berne}, although not explicitly limited to, nor directly
referring to folklore, member Union countries may use their respective national
legislation to designate a competent authority to represent, protect and enforce
the rights of an unknown, but assumed to be a domiciled in said country, author
of an unpublished work. After which, that member of the Union must by way of
written declaration notify the WIPO Director General of said details, who in
turn will communicate the details of the declaration to all other countries of
the Union, as per the introductory observations of
\cite{wipo85_model_provi_national_laws_folklore}.\\

In 1980 a working group was established between WIPO and UNESCO, where a study
for the draft model provisions for national laws on the protection of creations
of folklore were undertaken. In 1982 the Directors General of UNESCO and WIPO
convened a Committee of Governmental Experts on the Intellectual Property
Aspects of the Protection of Expressions of Folklore, where the \emph{`Model Provisions for National Laws on the Protection of Expressions of Folklore
Against Illicit Exploitation and Other Prejudicial Actions}' were adopted.
Moreover these provisions were designed in such a way as to be adaptable to
appropriate national conditions, with the primary objective of protecting
against the abuse [section 1]\cite{wipo85_model_provi_national_laws_folklore}, of
expressions of folklore whilst also encouraging it's further development and
dissemination, as per the introductory observations of
\cite{wipo85_model_provi_national_laws_folklore}.\\

[Section 2]\cite{wipo85_model_provi_national_laws_folklore} articulates the
expressions of folklore protected under these Model Provisions, as
\emph{``productions consisting of characteristic elements of the traditional artistic
heritage developed and maintain by a community''}, including verbal, musical,
physically demonstrative and tangible expressions. [Sections 3 and
4]\cite{wipo85_model_provi_national_laws_folklore} deal with the authorized
utilization and corresponding exceptions respectively. Permission from a
competent authority is required for: any publication, reproduction or any
distribution of copies of [section
3(i)]\cite{wipo85_model_provi_national_laws_folklore}; or any public recitation,
performance or transmission [section
3(ii)]\cite{wipo85_model_provi_national_laws_folklore} of expressions of folklore.
With the exception of utilization that is educational [section
4.1(i)]\cite{wipo85_model_provi_national_laws_folklore}; an illustration
consistent with fair practice [section
4.1(ii)]\cite{wipo85_model_provi_national_laws_folklore}; borrowed and re-used in
the creation of an original work of an author(s) [section
4.1(iii)]\cite{wipo85_model_provi_national_laws_folklore}; or where the
utilization is incidental [section
4.2]\cite{wipo85_model_provi_national_laws_folklore}.\\

Utilization of an expression of folklore without either consent from a competent
authority nor appropriate acknowledgement of source of the author or community
and/or geographic place from which the work originates [section
5.1]\cite{wipo85_model_provi_national_laws_folklore}, constitute offensive
violations [section 6]\cite{wipo85_model_provi_national_laws_folklore} which are
liable to seizure or other actions [section
7]\cite{wipo85_model_provi_national_laws_folklore} and other civil remedies such
as damages [section 8]\cite{wipo85_model_provi_national_laws_folklore}. Where the
definition of a `competent or supervisory authority' [section
9]\cite{wipo85_model_provi_national_laws_folklore} is left open to the relevant
member Union country.\\

[Section 10[\cite{wipo85_model_provi_national_laws_folklore} details the necessary
procedures required for authorization with respect to utilization of an
expression of folklore. Such requests should be in writing to the competent
authority / community concerned [section
10.1]\cite{wipo85_model_provi_national_laws_folklore}. Where a fee may be
collected, as established by the supervisory authority, for the purposes of
promoting / safeguarding national culture and folklore [section
10.2]\cite{wipo85_model_provi_national_laws_folklore}. The Model Provisions also
accommodate appeals from authorization applicant against decisions made by the
competent authority [section
10.3]\cite{wipo85_model_provi_national_laws_folklore}.\\

As per the provisions of [section
11]\cite{wipo85_model_provi_national_laws_folklore}, with respect to matters
relating to appeals or offenses, the question of jurisdiction is left open and
to be defined by the relevant member Union country. This vague and opened-ended
definition in the Model Provisions is indeed problematic and inadequate,
especially in light of electronic transactions and the nature of potential
infringements and offenses that span multiple jurisdictions, across a number of
member Union countries. As per the provisions of [section
12]\cite{wipo85_model_provi_national_laws_folklore} the Model Provisions shall in
no way limit, prejudice nor supersede any protections applicable to expressions
of folklore, enjoyed under copyright law; industrial property law; laws
protecting producers, performers and broadcasting organizations; or any other
laws or international treaties to which the country is a party to. Lastly
[sections 13 and 14]\cite{wipo85_model_provi_national_laws_folklore} describe how
the Model Provisions cannot be interpreted in any way to hinder the use or
development of expressions of folklore, and where expressions of folklore are
developed in a foreign country they are subject to reciprocal or international
treaty based protection.\\

While it is apparent that the Model Provisions provide a robust framework
enabling Union member countries to enforce protection of their communities'
expressions of folklore with the freedom and flexibility to be adapted to locale
specific concessions. There however exists a clear need for the Model Provisions
to be adapted and updated to cope with the problems facing developing nations
with respect to traditional knowledge and expressions of folklore associated
with the digital age.

\section{Identify the issues relating to the protection of ethno-botanical knowledge. In your answer refer to the following [25]:}
\label{sec:org19732ea}

\subsection{Protection against exploitation of industrial property. (10)}
\label{sec:org498c0d3}

With respect to the protection of an indigenous peoples ethno-botanical traditional knowledge, a number of inter-governmental organisational
negotiations with respect to amendments to the TRIPS \cite{wto17_trips} Agreement
on Traditional Knowledge have unfortunately collapsed,
\cite{saurombe09_protection_indigenous_trad_knowl_through_ip}, and the current
protection afforded is through [article 27(3)(b)]\cite{wto17_trips}, where member
states are empowered to consider protection of traditional knowledge through
intellectual property systems.\\

With respect to protection against exploitation of intellectual property, at the
third session of the Standing Committee on the Law of Patents of WIPO, it was
unsuccessfully proposed that a notification requirement be established. In that
where the subject matter of a patent application is based on the genetic
resources forming part of the ethno-botanical heritage of an indigenous peoples,
then a copy of the contract affording access to the genetic resources in their
country of origin should be filed, requiring possible amendments to \cite{wipo17_pct_ai,wipo70_pct}.\\

Alternatively, were an indigenous peoples endeavors to have their traditional
ethno-botanical knowledge documented and published as identifiable and
searchable prior art, then this would bring into dispute the novelty requirement
of any inventions based on said knowledge. Moreover, this would provide a
mechanism with which said indigenous peoples could pursue revocations of
registered patents granted to the predatory enterprises, albeit at substantial
legal costs. Or perhaps, where such patents are granted, provide co-ownership to
original inventors / indigenous peoples
\cite{amechi15_leverag_trad_knowl_medicinal}.\\


\subsection{Protection for exploitation of industrial property. (15)}
\label{sec:orgbd78049}

\subsubsection{Patent Law}
\label{sec:org7f86df6}

Indigenous peoples as the rights holders of traditional ethno-botanical
knowledge may protect such knowledge under the patent system as per the
provisions of [section 25(1)]\cite{rsa78_patents_act} within the Republic.
Provided however that the ethno-botanical knowledge in question is not for any
variety of animal or plant or any essentially biological process for the
production of animals of plants, not being a micro-biological process or the
product of such processes as per [section 25(4)(b)]\cite{rsa78_patents_act}. This
industrial property system however, presents a significant basic problem in that
emphasis is given to the `active ingredients' which have undergone expensive
isolation and testing \cite{simelane09_african_trad_knowl_sys_biodiv_manag}.

\subsubsection{Unfair Competition in respect of trade secrets}
\label{sec:orgf8c51a3}

Indigenous peoples must show that the information relating to their traditional
ethno-botanical knowledge is:
\begin{itemize}
\item \textbf{Confidential:} The knowledge \uline{discovered} by the indigenous peoples
constitutes information that was neither public knowledge nor released in the
public domain.
\item Has \textbf{trade value}: There exists a \uline{market} interested in the traditional
ethno-botanical knowledge of the indigenous peoples.
\end{itemize}

There are a number of examples in South African case law that have qualified for
the protection of secret information, specifically in regards to a manufacturing
process \emph{for example Harvey Tiling Co (Pty) Ltd v Rodomac (Pty) Ltd \& Another
1977 (1) SA 316 (T)}. Another example that was successfully tried as a matter
for the protection of secret information, was
\citetitle{dijkhorst81_atlas_v_pikkewyn}, \cite{dijkhorst81_atlas_v_pikkewyn}, where
judge Dijkhorst was famously quoted that this branch of law was designed to
\emph{``address the schemes of geniuses bent upon reaping what they have not sown.''}
One last successful example which qualified for the protection of secret
information was \citetitle{diemont72_stellwinetrust_v_oudemeester}
\cite{diemont72_stellwinetrust_v_oudemeester}, where Judge Diemont stipulated that
one is acting unlawfully and dishonestly when he `filches' information from a
competitor, devised through the skill and industry of that competitor for his
own profit.\\

Although \cite{wipo83_paris_conve_protect_ip} does not provide and substantive
treatment for secret information with regards to unfair competition, pertinent
provisions are accommodated within [article 39(1)]\cite{wto17_trips} through the
stipulation [article 10\(^{\text{textnormal\{bis\}}}\)]\cite{wipo83_paris_conve_protect_ip}.
Moreover [article 39(2)]\cite{wto17_trips} provides the framework to prevent
unlawful competition through the use of secret information by, either being used
by, disclosed to or acquired by others without their consent in a manner
contrary to honest practices, provided:
\begin{itemize}
\item The information is \textbf{indeed secret and not generally know}, as per [article
6(i)]\cite{wipo96_model_provi_unfair_comp} and [article
39(2)(a)]\cite{wto17_trips}.
\item The information must have \textbf{commercial value}, as per [article
6(ii)]\cite{wipo96_model_provi_unfair_comp} and [article
39(2)(b)]\cite{wto17_trips}.
\item Lastly, as per [article 6(iii)]\cite{wipo96_model_provi_unfair_comp} [article
39(2)(c)]\cite{wto17_trips}, the applicant was required to \textbf{keep the information
secret} subject to reasonable circumstances.
\end{itemize}

Within the context of secret information as it pertains to unlawful competition,
the internation instruments corroborate South African case law, for the most
part. Where they may be perceived to perhaps deviate very slightly, is that the
intentional instruments set the onus on the plaintiff, which in this instance
would be the indigenous people, to ensure that he has provided adequate
safeguards in ensuring that their secret information is indeed kept secret.

\subsubsection{Transfer of technology, licenses  and compensatory liability}
\label{sec:org531a7f0}

Lastly, it'll be briefly mentioned that indigenous people may as per the
provisions of [section 53]\cite{rsa78_patents_act} within the Republic, or those
of [article 28.2]\cite{wto17_trips} assign, transfer by success and conclude
licensing contracts with respect of their traditional ethno-botanical knowledge.
Where as per the provisions of [section 65]\cite{rsa78_patents_act}, they could
seek relief in the form of an interdict or damages against an enterprise
exploiting their traditional ethno-botanical knowledge.

\printbibliography
\end{document}